\documentclass[10pt,twocolumn]{article}

\usepackage[margin=0.85in]{geometry}
\usepackage{amsmath,amssymb}
\usepackage{booktabs}
\usepackage{microtype}
\usepackage{xcolor}
\usepackage{tikz}
\usepackage{pgfplots}
\pgfplotsset{compat=1.16}
\usepackage[hidelinks]{hyperref}
\usepackage{caption}
\captionsetup{font=small,labelfont=bf}

\definecolor{pgood}{RGB}{31,119,80}
\definecolor{pbad}{RGB}{200,40,40}
\definecolor{pnode}{RGB}{38,70,120}
\definecolor{pmid}{RGB}{210,140,20}
\definecolor{pgray}{RGB}{110,110,110}
\newcommand{\code}[1]{\texttt{\small #1}}
\newcommand{\todo}[1]{{\color{pbad}\textbf{[TODO: #1]}}}

\title{\vspace{-3.5em}\bfseries pillar-UDP Redundancy: Header Format and the
Loss/Bandwidth Trade\\[0.3em]\large\normalfont(companion to ``Clustered Multipath
Congestion Avoidance and Redundant Delivery'')}
\author{\normalsize pillar engineering}
\date{\normalsize DRAFT --- \today}

\begin{document}
\maketitle

\begin{abstract}
\noindent
\textbf{This is a working draft.} The primary pillar-UDP paper proves the
\emph{qualitative} delivery guarantees (exactly-once, route-around liveness,
bounded fan-out) and defers two quantitative matters to this companion: (i)~the
on-the-wire encoding of the per-connection \emph{redundancy allowance} and the
rest of the pillar-UDP message header, and (ii)~the analytic loss/bandwidth
trade that determines how large the allowance should be for a target
reliability. This draft fixes the analytic framework---which is closed-form and
stated here in full---and lays out the header's logical field set; the exact
byte encoding and the empirical validation of the independence assumption remain
open (\S\ref{sec:open}).
\end{abstract}

\section{Scope}
The companion's job is narrow: give operators a defensible rule for setting the
redundancy allowance, and give implementers a stable header contract. Nothing
here changes the protocol's proven guarantees; it parameterises them. Two
regimes share one substrate and differ only in redundancy shape:
\emph{small-message} delivery (a message rides in $W$ whole redundant copies)
and \emph{bulk} delivery (a payload is split into $K$ data shards plus $M$
erasure shards and reconstructed from any $K$ arrivals). Both are analysed
below.

\section{The loss/bandwidth trade (small messages)}\label{sec:trade}
Let a message be sent as $W$ redundant copies over paths whose losses are
independent with per-path probability $p$ (the primary paper argues dispersion
across AS-diverse nodes is what makes independence realistic; \S\ref{sec:open}
records this as the assumption to validate empirically). The message is
undelivered only if every copy drops:
\[
  P_{\text{fail}}(W,p) = p^{\,W}.
\]
The bandwidth cost is a fan-out of $W$: overhead is $(W-1)$ redundant copies per
delivered message. Fixing a target residual $\varepsilon$ and solving
$p^{W}\le\varepsilon$ gives the \emph{minimum allowance}
\[
  W^\star(p,\varepsilon)=\left\lceil \frac{\ln\varepsilon}{\ln p}\right\rceil
    = \left\lceil \log_{p}\varepsilon \right\rceil ,
\]
a closed form the sender can evaluate per connection from its measured $p$ and
the connection's declared $\varepsilon$. This is the value the redundancy
allowance encodes; it scales \emph{down} toward $W=1$ as $p\to0$ (a clean path
costs nothing) and \emph{up} as $p\to1$, which is exactly the adaptive posture
the primary paper requires. Table~\ref{tab:wstar} tabulates $W^\star$.

\begin{table}[t]
\centering\small
\begin{tabular}{@{}lrrr@{}}
\toprule
$p$ & $\varepsilon{=}10^{-3}$ & $\varepsilon{=}10^{-6}$ & $\varepsilon{=}10^{-9}$\\
\midrule
0.01 & 2 & 3 & 5 \\
0.05 & 3 & 5 & 7 \\
0.10 & 3 & 6 & 9 \\
0.20 & 5 & 9 & 13 \\
0.30 & 6 & 12 & 18 \\
0.50 & 10 & 20 & 30 \\
\bottomrule
\end{tabular}
\caption{Minimum redundancy allowance $W^\star=\lceil\log_p\varepsilon\rceil$ for
target residual loss $\varepsilon$. The allowance is cheap on good paths and
grows only as the path degrades.}
\label{tab:wstar}
\end{table}

\section{Erasure regime (bulk)}\label{sec:erasure}
For bulk transfer the payload is split into $K$ data shards and $M$ parity
shards ($N=K+M$ total, systematic code), reconstructable from any $K$ arrivals
(the implementation's \code{encode}/\code{reconstruct} in \code{pillar\_udp.rs}).
With independent per-shard loss $p$, reconstruction fails iff fewer than $K$ of
$N$ shards arrive:
\[
  P_{\text{fail}}(K,M,p)=\sum_{i=0}^{K-1}\binom{N}{i}(1-p)^{i}p^{\,N-i}.
\]
The bandwidth overhead is $M/K$, tunable independently of $K$, so bulk pays a
\emph{fractional} overhead where small messages pay a multiple. The parity rate
$M/K$ is tied to the same redundancy-allowance header field, mapped through the
$\varepsilon$/$p$ relation above rather than a raw copy count.
\todo{plot $P_{\text{fail}}(K,M,p)$ surface and derive the $M/K$ that meets a
given $\varepsilon$ at measured $p$; cross-check against Table~\ref{tab:wstar}.}

\section{Header format (logical fields)}\label{sec:header}
The header binds every datagram---initial spray copy or forward---to its logical
message and states the connection's redundancy contract up front. The
\emph{logical} field set below is fixed by the protocol's proven mechanics; the
concrete byte layout is deferred (\S\ref{sec:open}).

\begin{table}[t]
\centering\small
\begin{tabular}{@{}p{0.27\columnwidth}p{0.63\columnwidth}@{}}
\toprule
field & meaning \\
\midrule
CID & content address of the message bytes; the dedup and integrity key
(identical to the streamdb/blob content address). \\
redundancy allowance & the connection's $W^\star$ (small) or parity rate $M/K$
(bulk), per \S\ref{sec:trade}--\ref{sec:erasure}; caps system-wide fan-out. \\
hop TTL & remaining forward budget; decremented per hop, forwarding stops at 0
(with CID dedup, guarantees loop-free termination). \\
side-effect class & idempotent/\emph{Convergent} vs \emph{Exclusive}; routes an
exclusive message's processing to the coordination-core lease holder. \\
shard descriptor & for bulk: $K$, $M$, and the shard index/offset for
reconstruction. \\
return-routability token & anti-amplification: proof the source address is
reachable before the cell commits redundant replies. \\
\bottomrule
\end{tabular}
\caption{pillar-UDP header logical fields. Each corresponds to an invariant
already modelled in \code{specs/PillarUDP.tla} / \code{PillarUdpClient.tla} /
\code{PillarUdpMesh.tla}; this paper's remaining work is their wire encoding.}
\label{tab:header}
\end{table}

\section{Open items}\label{sec:open}
This draft is complete on the analytic framework and the logical header
contract; the following remain before it is publishable:
\begin{itemize}\itemsep2pt
  \item \todo{Byte-level wire encoding of Table~\ref{tab:header}: field widths,
  varint vs fixed, alignment, and the Noise/session-key framing boundary.}
  \item \todo{Empirical validation of the independence assumption
  (\S\ref{sec:trade}): measure correlated loss across AS-diverse pillar-ipam
  prefixes and quantify how much it inflates the effective $p$ used in
  $W^\star$.}
  \item \todo{Bulk $M/K$ selection curve (\S\ref{sec:erasure}) and its mapping
  onto the single redundancy-allowance field shared with the small-message
  regime.}
  \item \todo{Congestion-safety envelope: bound the aggregate fan-out a cell
  emits so the adaptive allowance cannot itself induce congestion collapse under
  correlated loss.}
\end{itemize}
No quantitative claim in this draft rests on unmeasured data: Tables~\ref{tab:wstar}
is exact arithmetic from the closed forms, and every empirical question is listed
above rather than filled with a placeholder value.

\end{document}
